\documentclass[11pt]{article}

\usepackage[margin=1.1in]{geometry}
\usepackage[T1]{fontenc}
\usepackage[utf8]{inputenc}
\usepackage{lmodern}
\usepackage{microtype}
\usepackage{booktabs}
\usepackage{amsmath}
\usepackage{graphicx}
\usepackage{xcolor}
\usepackage{listings}
\usepackage{multirow}
\usepackage[hidelinks]{hyperref}
\usepackage{caption}
\captionsetup{font=small,labelfont=bf}

\lstdefinestyle{code}{
  basicstyle=\ttfamily\small,
  breaklines=true,
  frame=single,
  framerule=0.3pt,
  xleftmargin=1em,
  aboveskip=0.8em,
  belowskip=0.8em,
  columns=fullflexible,
  keepspaces=true
}

\newcommand{\jssam}{JS-SAM}
\newcommand{\tlaplus}{TLA\texorpdfstring{\textsuperscript{+}}{+}}

\title{\textbf{Executable JavaScript as a Specification Language:\\
A \jssam{} Case Study on SysMoBench}}

\author{
  Jean-Jacques Dubray\\
  \texttt{jdubray@gmail.com}\\[0.5em]
  \small\textit{In collaboration with the SysMoBench / Specula team}
}
\date{July 2, 2026 \quad (draft)}

\begin{document}
\maketitle

\begin{abstract}
SysMoBench evaluates whether large language models (LLMs) can formally model
real-world concurrent and distributed systems, grading generated \tlaplus{}
specifications through four automated phases: syntax, bounded model checking,
conformance against execution traces captured from the real system, and
invariant verification. We extend SysMoBench with its first
\emph{non-formal-language} backend, \jssam{}: specifications written in
JavaScript following the SAM pattern, whose semantics --- actions propose,
a model accepts, state is a pure function of the model --- are deliberately
aligned with \tlaplus{}. The backend tests a specific hypothesis: LLMs may
model systems more accurately in a language ubiquitous in their training data
than in formal languages they rarely see.

Using the Asterinas OS spinlock task, we instrument the real kernel, run it
under QEMU, and capture a 28-window transition corpus; we then evaluate four
Claude models (Haiku~4.5, Sonnet~4.6, Opus~4.8, Fable~5) end to end. Four
findings emerge. First, trace conformance is the only phase with discriminating
power: all four models pass syntax, model checking, and invariant verification
(one of them, a metric audit later showed, vacuously --- its exploration
never left the initial state), while conformance spreads them from 21.4\% to 89.3\% --- including a frontier
model that silently no-ops an entire action (every \texttt{ReleaseLock}) yet
passes every internal-consistency phase. Second, a single round of
repair-from-trace-feedback splits the models into repairers and a
non-repairer (a replication showed the apparent middle ``partial'' tier does
not survive resampling), and shows that one-shot accuracy and repair ability
are different axes: the model with the categorical release defect repairs to
100\% --- with a generalization audit confirming the repairs are root-cause
fixes, not memorized windows --- while the weakest model cannot use the
feedback at all.
Third, on the same task and models, 4/4 \jssam{} specifications are
end-to-end checkable versus 1/4 for \tlaplus{}, attributable to a concrete
\tlaplus{} semantic pitfall (unbounded \texttt{CHOOSE}) that JavaScript's
native \texttt{null} sidesteps, and to \jssam{}'s self-contained executable
artifact eliminating the model-checker configuration surface. Finally, a
pre-registered, paired, four-arm study ($N{=}5$ generations per cell) puts
the conformance question itself under control. Most of an apparent \tlaplus{}
conformance advantage was \emph{prompt prescriptiveness}, not language:
adding the identical semantics block to the JavaScript prompt closed most of
the gap. A one-directional residual survives that control --- \tlaplus{}
reaches 100\% conformance for every model and every generation, constrained
\jssam{} does so for one model of four --- but a plain-JavaScript control arm
locates its cause: a bare \texttt{next(state, action, data)} transition
function, structurally isomorphic to the \tlaplus{} relation, ties \tlaplus{}
at 100\% for every model and every generation. The residual is therefore
attributable neither to JavaScript nor to formality but to the SAM
specification contract's defect surface (proposal/acceptor wiring plus
mandatory auxiliary state); what models transcribe reliably is
\emph{minimal declarative shape}, in any language. The hypothesis splits:
familiarity and a self-contained artifact buy checkability; transcription
fidelity is governed by the shape of the specification contract --- a
directly actionable design result for \jssam{}, which we validate at
$N{=}5$ on two tasks --- the kernel spinlock and a distributed lock
service --- where 40 of 40 model-generated lean-contract specifications
pass every phase. We report these results
candidly as a single-task case study, describe the kernel trace-capture and
controlled-comparison methodology in full, and outline the path to a
multi-task comparison.
\end{abstract}

\section{Introduction}
\label{sec:intro}

Formal specifications are among the most effective tools for establishing the
correctness of concurrent and distributed systems~\cite{newcombe2015amazon},
and among the least used: writing them demands rare expertise and sustained
effort. LLMs raise an obvious question --- can they write specifications for
us? --- and a harder one: \emph{how would we know if they got it right?}
A generated specification can be syntactically flawless, explore cleanly under
a model checker, satisfy plausible invariants, and still describe a system
that does not exist. The SysMoBench benchmark~\cite{sysmobench} confronts this
directly. Its authors observed that when asked to model Etcd's Raft
implementation, a frontier LLM produced, in effect, the specification from the
appendix of the Raft paper --- a ``textbook'' model with little connection to
the implementation at hand~\cite{sigops2026}. SysMoBench therefore grades
specifications in four phases of increasing demand: (1)~syntax, (2)~runtime
model checking, (3)~\emph{transition validation} --- replaying
$(\mathit{pre}, \mathit{action}, \mathit{post})$ windows captured from real
executions against the generated model --- and (4)~invariant verification.

SysMoBench was built for \tlaplus{}~\cite{lamport2002specifying}, and has since
grown a pluggable backend interface with Alloy~\cite{jackson2002alloy} and
PAT/CSP\#~\cite{sun2009pat} backends. All three are formal languages, and all
three are rare in LLM training corpora relative to mainstream programming
languages. This asymmetry motivates the extension we present here.

\paragraph{The \jssam{} hypothesis.}
SAM~\cite{sam,dubray2016infoq} is a software-engineering pattern whose
semantics are explicitly derived from \tlaplus{}: \emph{actions} compute
proposals from inputs, a \emph{model} accepts or rejects proposals as the sole
locus of mutation, and \emph{state} is a pure function of the model. A SAM
specification is an ordinary executable JavaScript module. JavaScript is
plausibly the single most abundant programming language in LLM training data.
The hypothesis, then:

\begin{quote}
\emph{Can an LLM produce a correct specification of a real system more readily
in a pattern hosted by a language it has seen constantly in training
(JavaScript), than in formal languages it has rarely seen --- while preserving
the \tlaplus{}-style discipline that makes the artifact checkable?}
\end{quote}

Either answer is informative. If \jssam{} scores diverge upward, host-language
familiarity dominates and formal syntax is a bottleneck; if scores match or
fall, the hard part is modeling, not notation.

\paragraph{Contributions.}
This paper reports the first end-to-end results from the \jssam{} backend,
all on the SysMoBench \texttt{spin} task (the Asterinas OS
spinlock~\cite{asterinas}):

\begin{enumerate}
\item \textbf{A non-formal-language backend for SysMoBench}
  (\S\ref{sec:backend}): the \jssam{} module contract, a sandboxed execution
  helper, and --- a first for the benchmark --- a \emph{direct} Phase-3
  replay path requiring no coding-agent orchestration, because SAM's step
  relation \emph{is} $\mathit{model.present}(\mathit{action}(\mathit{data}))$.
\item \textbf{A reproducible kernel trace-capture methodology}
  (\S\ref{sec:traces}): instrumenting the real Asterinas spinlock, driving a
  two-thread contention scenario in a single-CPU kernel test without deadlock,
  and folding low-level kernel events into model-independent
  $(\mathit{pre}, \mathit{action}, \mathit{post})$ windows --- producing the
  28-window corpus used throughout.
\item \textbf{A four-model comparison} (\S\ref{sec:models}) showing that
  Phases 1, 2, and 4 have no discriminating power on this task while Phase 3
  spreads the models from 21.4\% to 89.3\% --- and that general model
  capability does not predict modeling accuracy: the most capable model tested
  produced a specification that no-ops every lock release.
\item \textbf{A repair-loop experiment} (\S\ref{sec:repair}) demonstrating
  that self-correction from precise trace feedback is a distinct capability
  axis from one-shot accuracy --- capable models repair fully (verified by a
  held-out generalization audit as fixes, not memorization), the weakest not
  at all, with mid-tier outcomes unstable at $N{=}1$.
\item \textbf{A \jssam{}-vs-\tlaplus{} head-to-head} (\S\ref{sec:vstla}):
  with identical models and task, 4/4 \jssam{} specifications reach
  end-to-end evaluability versus 1/4 for \tlaplus{}, with each \tlaplus{}
  failure attributed to a specific structural cause.
\item \textbf{A controlled cross-language conformance study}
  (\S\ref{sec:crosslang}): a pre-registered, paired, four-arm design
  ($N{=}5$ generations per cell) that pins both languages to the same
  observable-state contract and replays the same trace corpus mechanically in
  both. After catching and controlling a prompt-prescriptiveness confound
  that would otherwise have produced a false headline, a plain-JavaScript
  control arm attributes the residual conformance gap to the SAM contract's
  machinery --- not to the JavaScript language, and not to formality ---
  yielding a concrete design change for the backend, validated by 40/40
  lean specifications passing every phase across two tasks
  (\S\ref{sec:locksvc}).
\end{enumerate}

We emphasize scope honestly: the controlled experiments concern one system
--- the spinlock --- and a 28-window trace corpus (the lean-contract
validation extends to a second task with its own 60-window corpus);
Experiments 1--3 are single-generation, with $N{=}5$ replication in
Experiment 4 and the lean studies. We present it as a case study whose
value lies in the precision of the failure attribution and the reproducibility
of the pipeline, not as a settled verdict on the hypothesis
(\S\ref{sec:limitations}).

\section{Background}
\label{sec:background}

\subsection{SysMoBench}

SysMoBench~\cite{sysmobench} pairs eleven real systems (concurrent
synchronization primitives and distributed protocols, from the Asterinas
spinlock to Etcd Raft) with instrumentation harnesses and expert-written
invariants. Given the system's source, an LLM generates a specification, which
is scored in four automated phases:

\begin{description}
\item[Phase 1 --- Syntax.] Does the specification parse and structurally
  validate? (For \tlaplus{}: the SANY analyzer.)
\item[Phase 2 --- Runtime model checking.] Does bounded state-space
  exploration terminate without errors, deadlocks, or nondeterminism? (For
  \tlaplus{}: TLC~\cite{yu1999tlc}.)
\item[Phase 3 --- Transition validation (conformance).] Execution traces from
  the \emph{real} system are cut into windows
  $(\mathit{pre}, \mathit{action}, \mathit{post})$; each window passes if the
  generated model, started at $\mathit{pre}$ and applying $\mathit{action}$,
  produces $\mathit{post}$. The output is a per-action scorecard.
\item[Phase 4 --- Invariant verification.] Expert invariants (e.g., mutual
  exclusion) are translated into the specification's language and checked.
\end{description}

The benchmark's central published finding is that frontier LLMs cluster near
100\% on syntax but average roughly 46\% on conformance and 41\% on
invariants~\cite{sigops2026}: models write valid \tlaplus{} fluently yet
describe the textbook protocol rather than the implementation. Phase 3 exists
precisely to distinguish ``writing \tlaplus{}'' from ``modeling this system.''

\subsection{The SAM pattern}

SAM (State--Action--Model)~\cite{sam,dubray2016infoq} factors a program's
temporal behavior the way \tlaplus{} factors a specification. An
\emph{action} is a pure function from inputs to a \emph{proposal}; the
\emph{model} alone decides whether to accept a proposal and mutate itself
(the analogue of a \tlaplus{} next-state relation, with acceptors playing the
role of guards); \emph{state} is a pure function of the model. A step of the
system is $\mathit{model.present}(\mathit{action}(\mathit{data}))$. The
pattern is implemented by the \texttt{sam-pattern}
library~\cite{sampattern}, which also provides a bounded explorer
(\texttt{checker}) over the model's reachable states; in the \jssam{}
backend, this checker is what executes Phases 2 and 4.

Two properties make SAM suitable as a specification substrate. First, the
semantic alignment with \tlaplus{} means the benchmark's phases map onto it
without distortion: bounded exploration is Phase 2, replaying a trace window
is literally one SAM step, and invariants are predicates over
$\mathit{getState}()$. Second, the artifact is an executable JavaScript
module: there is no separate tool chain between the specification and its
own evaluation.

\section{The \jssam{} backend}
\label{sec:backend}

SysMoBench backends implement a \texttt{LanguageBackend} interface and
register themselves; the CLI and the four evaluators resolve them by name with
no further plumbing. The \jssam{} backend consists of a Python adapter and a
small Node.js helper (\texttt{tools/js-sam/cli.mjs}) with four subcommands ---
\texttt{validate}, \texttt{check}, \texttt{transitions},
\texttt{invariants} --- one per phase.

\paragraph{Module contract.}
The generation prompt instructs the model to emit a single CommonJS module:

\begin{lstlisting}[style=code]
module.exports = {
  instance,       // SAM instance (synchronous steps)
  init,           // () => void : reset to the initial state
  actions,        // { AcquireLock: (data)=>void, ReleaseLock: (data)=>void }
  getState,       // () => plain JSON-serializable snapshot
  setState,       // (snapshot) => void : force a state (Phase 3 replay)
  checkerIntents, // action descriptors + input domains (Phases 2 & 4)
};
\end{lstlisting}

\noindent
\texttt{actions} must cover exactly the task's target actions;
\texttt{getState} must return a plain object (this is what trace states are
diffed against); \texttt{setState} exists solely so Phase 3 can replay a
window without searching for a path to its pre-state; and the module must be
deterministic and I/O-free. \texttt{checkerIntents} declares each action's
finite input domain --- for \texttt{spin}, every
$\{\mathit{thread}, \mathit{callType}\}$ combination --- playing the role of a
\tlaplus{} constants block. Notably, there is \emph{no separate configuration
artifact}: the module is self-contained, a fact that becomes material in
\S\ref{sec:vstla}.

\paragraph{Sandboxing.}
Model-generated JavaScript is untrusted. Every helper invocation runs in a
throwaway \texttt{node:20-slim} Docker container with no network,
a read-only root filesystem, a non-root user, and CPU/memory/PID caps. The
container is the trust boundary.

\paragraph{The direct Phase-3 path.}
For \tlaplus{}, SysMoBench's transition validation drives TLC through a
coding-agent workflow that constructs a checking module per window. \jssam{}
needs none of that: the evaluator loads the module, and for each window
executes \texttt{setState(pre)}; \texttt{actions[name](data)}; and diffs
\texttt{getState()} against \emph{post} under the benchmark's projection rule
(only keys present in the trace's post-state must match, so model-internal
bookkeeping is ignored). This ``direct path'' is the first non-agent Phase-3
implementation in the benchmark, and it follows from SAM's semantics: a
window \emph{is} a SAM step.

\paragraph{Cross-platform engineering.}
Making the pipeline run on a Windows host (in addition to Linux/macOS)
required fixing several host assumptions in the harness --- POSIX-only
\texttt{os.getuid()} calls, host-path container mounts invalid for
\texttt{C:\textbackslash} paths, and console encoding --- plus threading the
helper's \texttt{stepsExplored} metric through a field previously hardcoded to
zero for non-\tlaplus{} backends. We also generalized the benchmark's
agent-based invariant translator (previously \tlaplus{}-only) into a shared,
language-neutral component that \jssam{} can use, although all experiments
below use the direct translator for parity across languages. These changes
are upstreamed as pull requests.

\section{Capturing ground truth: kernel traces for Phase 3}
\label{sec:traces}

Phase 3 is the only phase requiring data external to the model, and SysMoBench
does not ship trace corpora: traces are generated on demand by instrumenting
the real system. Additionally, the maintainers' own \texttt{spin} captures use
three threads while the \jssam{} prompt models two, so the backend required
its own corpus at the granularity it models. This section documents the
pipeline; it is the experimental setup's main methodological contribution and
is fully reproducible from a single script.

\subsection{Pipeline}

\begin{enumerate}
\item \textbf{Pin a compatible revision.} The reference instrumentation patch
  applies cleanly at Asterinas v0.16.0 (matching the toolchain in the
  \texttt{asterinas/asterinas:0.16.0} image), not at \texttt{main}.
\item \textbf{Instrument.} Apply the benchmark's reference patch, which adds
  a trace module to the kernel's spinlock emitting one serial JSON event per
  lock operation, plus a new two-thread kernel test (below).
\item \textbf{Build and run.} Inside the container: build the initramfs and
  run the instrumented kernel test under QEMU (software emulation), collecting
  the event stream from the serial port.
\item \textbf{Parse.} Fold the event stream into
  $(\mathit{pre}, \mathit{action}, \mathit{post})$ NDJSON windows.
\end{enumerate}

\subsection{A two-thread contention scenario without deadlock}

The reference spinlock kernel tests are single-actor: every operation is
attributed to one thread, so contention --- the behavior that makes a lock
interesting --- never appears in the trace. But a blocking \texttt{lock()}
from a second actor in a single-CPU kernel test would spin forever. The
scenario therefore expresses the second actor's contention through
non-blocking \texttt{try\_lock} calls that fail without blocking:

\begin{lstlisting}[style=code]
actor0.lock()      -> TryAcquireBlocking(0), AcquireSuccess(0)  // 0 holds
actor1.try_lock()  -> AcquireFail(1)          // contention, no deadlock
actor0.release()   -> Release(0)
actor1.lock()      -> TryAcquireBlocking(1), AcquireSuccess(1)  // 1 holds
actor0.try_lock()  -> AcquireFail(0)          // contention
actor1.release()   -> Release(1)
actor0.try_lock()  -> AcquireSuccess(0)       // try from free succeeds
actor0.release()   -> Release(0)
\end{lstlisting}

A second, longer scenario --- more alternation of the holder, more failed and
successful \texttt{try\_lock} attempts, and same-thread reacquisition ---
extends the corpus from 8 windows to \textbf{28 windows} (17
\texttt{AcquireLock}, 11 \texttt{ReleaseLock}), the corpus used in all
experiments below.

\subsection{Model-independent state windows}

The kernel emits low-level events (\texttt{TryAcquireBlocking},
\texttt{AcquireSuccess}, \texttt{AcquireFail}, \texttt{Release}). The parser
folds these into windows keyed on the \emph{objective} spinlock state ---
$\{\mathit{lockHeld}, \mathit{lockHolder}\}$, the real observable ownership
--- and relies on the projection rule to ignore any model's internal
bookkeeping (thread status, call type). A failed acquire under contention
produces a window whose post-state ownership is unchanged; a successful
acquire records the new holder together with whether it arrived via the
blocking (\texttt{lock}) or non-blocking (\texttt{try}) path. Keying the
corpus on observable state keeps the ground-truth oracle independent of any
model under test.

\section{Experiment 1: four models, four phases}
\label{sec:models}

\subsection{Setup}

Four Claude models spanning the capability range --- Haiku~4.5, Sonnet~4.6,
Opus~4.8, and Fable~5 (Anthropic's most capable model at the time of
writing) --- each generated one \jssam{} specification for \texttt{spin} via
a single prompt (\texttt{direct\_call}), and the same artifact was scored in
all four phases. Phase 2 and 4 exploration ran through the
\texttt{sam-pattern} library's bounded checker at the task's default depth
bound of 6; Phase 3 used the 28-window corpus; Phase 4 checked three expert
invariants (MutualExclusion, LockStatusConsistency, NoDeadlock) translated
via the direct API path.

\subsection{Results}

\begin{table}[ht]
\centering
\caption{Four-phase \jssam{} results on \texttt{spin} (28-window corpus).
P3 per-action columns give windows passed / windows of that action.}
\label{tab:models}
\small
\begin{tabular}{lccrccc}
\toprule
Model & P1 & P2 (distinct states) & \multicolumn{1}{c}{P3 overall} & P3 Acquire & P3 Release & P4 \\
\midrule
Claude Opus 4.8   & pass & pass (7) & \textbf{89.3\%} (25/28) & 14/17 & 11/11 & 3/3 \\
Claude Fable 5    & pass & pass (7) & \textbf{50.0\%} (14/28) & 14/17 & \textbf{0/11} & 3/3 \\
Claude Sonnet 4.6 & pass & pass (7) & \textbf{50.0\%} (14/28) & 14/17 & 0/11 & 3/3 \\
Claude Haiku 4.5  & pass & pass (\textbf{1} --- vacuous) & \textbf{21.4\%} (6/28)  & 6/17  & 0/11 & 3/3 \\
\bottomrule
\end{tabular}
\end{table}

Two results stand out.

\paragraph{Only the reality check discriminates.}
Every model passes Phase 1, Phase 2, and Phase 4. On this task those three
phases have \emph{zero} discriminating power. An earlier draft reported the
identical 326{,}592 ``states explored'' per model as confirming structurally
equivalent state spaces; an audit prompted by review showed that number is the
checker's \emph{step count} --- safety-callback invocations over the
intent-permutation tree, $(d{+}1)\cdot 6^{d} = 7\cdot 6^{6}$ --- a pure
function of the contract-pinned intent domain and depth bound, identical for
any conforming spec and therefore model-independent. Counting \emph{distinct
semantic states} (unique model snapshots) instead separates the specs: 7 for
Opus, Fable, and Sonnet, but \textbf{1 for Haiku}, whose intent actions drop
their arguments so every proposal is rejected --- its exploration never leaves
the initial state. Haiku's Phase-2 pass is thus \emph{vacuous} (no crash,
deterministic, and serializable hold trivially for a spec that cannot move),
and its Phase-4 invariants are satisfied over that single state. Phase 2
verifies checkability, not semantic adequacy; the distinct-state count, now
reported by the checker, is the semantic signal. Phase 3 spreads the models
across a $4\times$ range, and --- restricted to comparisons it can make ---
orders Opus above the mid tier above Haiku. This independently reproduces, in
a non-formal language, the central SysMoBench finding for \tlaplus{}: internal
consistency is cheap; conformance with the real system is where specifications
fail~\cite{sigops2026}.

\paragraph{General capability does not predict modeling accuracy.}
The most striking row is Fable~5, the most capable model tested, scoring
50\% --- below Opus~4.8 and tied with mid-tier Sonnet. The cause is specific
and systematic, not noise: Fable's specification handles acquisition
competently (14/17) but \emph{no-ops every single release} (0/11). After any
\texttt{ReleaseLock}, its model leaves the lock held:

\begin{lstlisting}[style=code]
ReleaseLock  expected {lockHeld: false, lockHolder: null}
             got      {lockHeld: true,  lockHolder: 0}
\end{lstlisting}

The release logic \emph{reads} correctly in isolation --- it guards on the
holder, then clears the state --- but on replay the release never takes
effect: a categorical, whole-action modeling defect. Crucially, this
specification passed syntax, completed the full bounded exploration (7
distinct states) without violation, and satisfied all three invariants,
including mutual exclusion. (Its release \emph{does} fire when reached from
the initial state, where its auxiliary \texttt{threadStatus} guard holds; the
defect is that replay pins only the observable pre-state, leaving the guard
unsatisfiable --- so exploration from init cannot expose it.) Only replay
against the real kernel's transitions caught it.

The recurring failure among the passing acquisitions is also precise: all
three 14/17 models fail exactly the three \texttt{try\_lock}-from-free
windows --- the non-blocking acquire that should succeed on a free lock.
Blocking acquisition (abundant in training data and in textbook treatments)
is modeled correctly by every model except Haiku; the less common
non-blocking path is where they slip. This echoes, at micro scale,
SysMoBench's ``textbook modeling'' diagnosis: models reproduce the familiar
template and miss the implementation-specific path.

\section{Experiment 2: repair from trace feedback}
\label{sec:repair}

\subsection{Setup}

If a model is shown \emph{exactly} which transitions its specification got
wrong --- the pre-state, the action and its data, the real system's
post-state, and what the model produced instead --- can it fix the
specification? For each model we ran one repair round: score the original
specification on Phase 3; build a repair prompt containing the full original
module plus every failing window; ask the model to rewrite the module;
confirm the result still passes Phase 1; re-score Phase 3 on the same
28-window corpus.

\subsection{Results}

\begin{table}[ht]
\centering
\caption{One-round repair from Phase-3 feedback (same 28-window corpus;
original run --- the generalization-audit replication differs for Sonnet,
which repaired to 28/28 there).}
\label{tab:repair}
\small
\begin{tabular}{lcccc}
\toprule
Model & Baseline P3 & Repaired P3 & Acquire (base$\to$rep) & Release (base$\to$rep) \\
\midrule
Claude Opus 4.8   & 89.3\% (25/28) & \textbf{100\%} (28/28) & 82\% $\to$ 100\% & 100\% $\to$ 100\% \\
Claude Fable 5    & 50.0\% (14/28) & \textbf{100\%} (28/28) & 82\% $\to$ 100\% & \textbf{0\% $\to$ 100\%} \\
Claude Sonnet 4.6 & 50.0\% (14/28) & 60.7\% (17/28) & 82\% $\to$ 100\% & 0\% $\to$ 0\% \\
Claude Haiku 4.5  & 21.4\% (6/28)  & 21.4\% (6/28)  & 35\% $\to$ 35\%  & 0\% $\to$ 0\% \\
\bottomrule
\end{tabular}
\end{table}

\paragraph{Self-correction separates the models --- but the middle tier is noise.}
In the run of Table~\ref{tab:repair}, Opus and Fable repair to a perfect
28/28 in a single round, Sonnet repairs \emph{partially} (it fixes the
\texttt{try\_lock} acquisition defect but not its release path), and Haiku
does not improve at all --- superficially a clean three-tier gradient
(full / partial / none). A replication run (part of the generalization
audit below, with artifacts saved) revised this: \emph{Sonnet also repaired
to 28/28}. The boundary that is stable across both runs is two-tier ---
Opus, Fable, and (unstably) Sonnet can exploit transition-level feedback;
Haiku cannot --- and single-round repair outcomes for mid-tier models are
sampling-dependent at $N{=}1$.

\paragraph{The Fable reversal.}
Fable's baseline looked no better than Sonnet's, and its failure mode ---
a whole action silently broken --- looked worse. Yet under repair, Fable
fully recovers: its release defect was a \emph{recoverable slip}, not a
capability ceiling. One-shot specification accuracy and
repair-from-feedback are different axes of capability, and a benchmark that
measures only the former will misrank models for any workflow that includes
a correction loop --- which realistic agentic formal-modeling workflows
do~\cite{sigops2026}. (An earlier draft contrasted Fable's recovery with
Sonnet's failure to fix the identical symptom; the replication run removed
that contrast.)

\paragraph{Generalization audit: the repairs are fixes, not memorization.}
Repaired specifications were originally re-checked only on Phases 1 and 3,
on the same 28 windows quoted in the repair prompt --- leaving open that a
repair might \emph{memorize} the failing windows (branch on the pre-state,
return the expected post), which would make this an experiment about
prompt-following rather than modeling. An audit re-ran the repair loop with
artifacts saved and scored each baseline and repaired specification on the
seen corpus, on \emph{held-out} windows, and on Phases 2 and 4. The 28 seen
windows collapse to 8 distinct (pre-state, action, data) combinations; the
held-out set is the 10 combinations of the observable domain the corpus
never shows (acquire on a held lock, release by a non-holder, release of a
free lock). All repaired specifications pass all held-out windows, reach
the same 7 distinct semantic states as a correct specification under
bounded exploration (Haiku's unrepaired specification remains at 1), and
hold all invariants. Direct inspection found no memorization --- no
state-equality tables or window literals; the successful repairs are the
same root-cause fix, moving the release ownership check from the auxiliary
\texttt{threadStatus} guard (unsatisfiable under a pinned observable
pre-state) to the authoritative \texttt{lockHeld}/\texttt{lockHolder},
with Sonnet's repair documenting the diagnosis in a comment. One
structural honesty note: because the kernel instrumentation emits acquire
events at acquisition success, every held-out combination is an observable
no-op, so held-out scoring alone refutes only memorizers with unsafe
defaults --- the inspection and Phase-2/4 evidence carry the conclusion,
and a state-changing held-out set (as \texttt{locksvc} affords) is the
right vehicle for a stronger version of this audit.

\paragraph{What this does and does not measure.}
The repair prompt contains the failing windows themselves, so the experiment
measures ``can the model exploit precise, correct feedback,'' not blind
improvement. That is the intended question --- it is exactly the signal
available inside an agentic repair loop --- but it is not a from-scratch
re-test, and we ran a single round per attempt; whether Haiku converges
with more rounds is untested.

\section{Experiment 3: \jssam{} versus \tlaplus{}}
\label{sec:vstla}

\subsection{Setup}

The head-to-head the backend exists to inform: the same four models, the same
\texttt{spin} task, one \texttt{direct\_call} generation per model per
language, scored on Phase 1 (syntax: SANY vs.\ module validation), Phase 2
(model checking: TLC vs.\ bounded exploration), and Phase 4 (invariant
verification, direct translator for both languages). Phase 3 is deliberately
\emph{not} compared: \jssam{}'s Phase 3 is a self-contained direct replay
while \tlaplus{}'s is a heavyweight coding-agent path --- different
mechanisms, not a fair one-to-one. Phase 4 is compared as pass/fail because
the two languages ship different invariant-template counts for \texttt{spin}
(seven for \tlaplus{}, three for \jssam{}).

\subsection{Results}

\begin{table}[ht]
\centering
\caption{\jssam{} vs.\ \tlaplus{} on \texttt{spin}: models producing a
passing artifact per phase.}
\label{tab:vstla}
\small
\begin{tabular}{lccc}
\toprule
 & P1 syntax & P2 model-checkable & P4 invariants \\
\midrule
\jssam{}  & 4/4 & \textbf{4/4} & \textbf{4/4} \\
\tlaplus{} & 4/4 & \textbf{1/4} & \textbf{1/4} \\
\bottomrule
\end{tabular}
\end{table}

\begin{table}[ht]
\centering
\caption{Per-model \tlaplus{} detail.}
\label{tab:tladetail}
\small
\begin{tabular}{lcccl}
\toprule
Model & P1 & P2 & P4 & P2 failure cause \\
\midrule
Claude Opus 4.8   & pass & \textbf{fail} & fail & unbounded \texttt{CHOOSE} (TLC rejects) \\
Claude Fable 5    & pass & pass (159 distinct states) & pass (7 inv.) & --- \\
Claude Sonnet 4.6 & pass & \textbf{fail} & fail & unbounded \texttt{CHOOSE} (TLC rejects) \\
Claude Haiku 4.5  & pass & \textbf{fail} & fail & config generation fell back; no usable \texttt{.cfg} \\
\bottomrule
\end{tabular}
\end{table}

\paragraph{Syntax is not where the languages differ.}
Every model writes syntactically valid specifications in \emph{both}
languages --- consistent with SysMoBench's finding that frontier models have
essentially saturated \tlaplus{} syntax~\cite{sigops2026}. The divergence is
at Phase 2: whether the specification is actually \emph{checkable}. All four
\jssam{} modules run; one of four \tlaplus{} specifications does. Each
\tlaplus{} failure is attributable to a specific structural cause.

\paragraph{Cause 1: a \tlaplus{} semantic pitfall JavaScript does not have.}
Opus and Sonnet both modeled ``no lock owner'' as
\texttt{NULL == CHOOSE v : v \textbackslash notin Threads}. SANY accepts
this --- syntax passes --- but TLC cannot evaluate an unbounded
\texttt{CHOOSE}, so model checking fails. The idiom is natural, looks
correct, and is a well-known TLC trap. JavaScript has a native \texttt{null},
so the corresponding \jssam{} specifications express the same concept with no
trap available to fall into. This structural hazard caught two of four
models, \emph{including the model with the best \jssam{} conformance score}
(Opus). It is a clean instance of the hypothesis's mechanism: the failure is
not ``the model cannot model a spinlock'' (its JavaScript model of the same
system scored 89.3\% against real traces) but ``the formal language exposes a
semantic surface on which competent models slip.''

\paragraph{Cause 2: a configuration surface \jssam{} does not have.}
\tlaplus{} model checking requires a separate \texttt{.cfg} artifact
(constants, init/next, invariants) that the harness must generate to match
the specification; for Haiku, that generation fell back and produced no
usable configuration (zero states explored). A \jssam{} module is
self-contained and executable --- there is no second artifact to generate,
so this entire failure mode is structurally absent. We flag the attribution
honestly: the \texttt{CHOOSE} failures belong to the models; the Haiku
failure belongs to the harness. But the asymmetry itself is structural ---
one language's evaluability depends on a fragile secondary artifact and the
other's does not.

\paragraph{Net.}
On this task, models reach an end-to-end evaluable specification far more
readily in \jssam{} (4/4 through Phase 4) than in \tlaplus{} (1/4). This is
direct --- though single-task --- evidence for the hypothesis in its
practical form: \emph{host-language familiarity plus a self-contained
executable artifact lowers the barrier to a usable specification.} It does
not, by itself, show that \jssam{} specifications are more \emph{faithful};
it shows they are far more often checkable at all, which is the precondition
for faithfulness to be measured. Experiment 4 (\S\ref{sec:crosslang})
measures faithfulness directly, under a design built to make the two
languages comparable --- and resolves it in an unexpected place.

\section{Experiment 4: cross-language conformance under a controlled contract}
\label{sec:crosslang}

Experiment 3 could not compare Phase 3 across languages because the two
backends implement it by incomparable mechanisms: \jssam{} runs a mechanical
replay, while \tlaplus{} runs a coding-agent path --- necessary because a
free-form \tlaplus{} specification invents its own variable names, so mapping
trace states onto spec states requires interpretation. A naive comparison
would measure the replay machinery, not the specifications. This experiment
removes that asymmetry with a pre-registered, paired design; the full
methodology note, written to be socialized with the benchmark maintainers, is
included in the artifact repository.

\subsection{Design}

\paragraph{Equalize the observable-state contract.}
The degree of freedom that forces \tlaplus{} onto the agent path --- free
choice of variables --- is one the \jssam{} prompt already removes by pinning
\texttt{getState()}'s shape. We remove it symmetrically: a constrained
\tlaplus{} prompt mandates \texttt{VARIABLES lockHeld, lockHolder} (the trace
schema), action operators \texttt{AcquireLock(thread, callType)} /
\texttt{ReleaseLock(thread)}, and a TLC-safe sentinel
(\texttt{NONE == "none"}, not \texttt{CHOOSE}). Both languages then answer
the identical question --- \emph{express this system over this observable
state} --- and trace states map onto spec states with no interpretation.

\paragraph{A direct TLC replay path (no agent), functional on both sides.}
With the contract fixed, \tlaplus{} transition validation becomes mechanical,
mirroring \jssam{}'s \texttt{setState} $\to$ action $\to$ diff: for each
window, a synthesized module pins the pre-state in its initial predicate and
takes exactly one step of the window's action. Mere reachability of the
traced post-state would be too weak a criterion: a \tlaplus{} action is a
\emph{relation}, and an over-permissive action can reach the correct
post-state while also admitting wrong ones --- passing an existential check
that the functional JS replay (one \texttt{next} state, compared exactly)
would fail. The replay therefore enforces the same functional relation on
both languages: it enumerates the action's one-step image from the pinned
pre-state and requires that the image be exactly the traced post-state
(branching factor 1; over-permissiveness fails the window). The check was
validated in three directions: a contract-following specification passes
28/28 (positive control); the same specification with \texttt{ReleaseLock}
mutated to a no-op fails exactly the 11 release windows (negative control);
and a deliberately permissive specification --- holder set to \emph{any}
thread on acquire --- passes 28/28 existentially but only 11/28
functionally, with branching factor 2 (permissive control), demonstrating
that the functional check is strictly stronger and necessary. An audit of
all 20 generated \tlaplus{} specifications found every one deterministic:
existential and functional scores coincide on every window, and the maximum
branching factor is 1. The \tlaplus{} results below are therefore earned
under a relation as strict as the JS one, not inflated by existential
slack.

\paragraph{Four arms, because the prompt and the contract are both confounds.}
An initial two-arm run (deployed \jssam{} prompt vs.\ constrained \tlaplus{})
produced what looked like a decisive \tlaplus{} reversal. But the constrained
\tlaplus{} prompt \emph{states the exact single-step observable semantics}
(free $\to$ acquire; held $\to$ unchanged; holder $\to$ release), which the
deployed \jssam{} prompt does not: a difference in prompt content, not in
language. And even with prompts equalized, the two substrates still differ in
a second way --- the SAM contract obliges an executable module with
proposal/acceptor wiring and auxiliary state, while the \tlaplus{}
specification is a bare two-variable relation. We therefore ran four arms:
\textbf{JS(deployed)} --- the shipped \jssam{} prompt;
\textbf{JS(constrained)} --- the deployed prompt plus the \emph{identical}
semantics block used by the \tlaplus{} arm;
\textbf{plain-JS} --- the same constrained semantics, but the specification
is a bare \texttt{next(state, action, data)} transition function over only
$\{\mathit{lockHeld}, \mathit{lockHolder}\}$: JavaScript, yet structurally
isomorphic to the \tlaplus{} relation, with no SAM library and no auxiliary
state; and \textbf{TLA(constrained)}. JS(constrained) vs.\ TLA(constrained)
varies language and contract together; plain-JS vs.\ TLA(constrained) varies
\emph{only} the language; JS(constrained) vs.\ plain-JS varies \emph{only}
the contract.

\paragraph{Replication, pairing, and checkability.}
$N{=}5$ generations per model per arm (80 in total), all scored per-window on
the same 28-window kernel corpus. The pre-registered analysis was a pooled
per-window McNemar test; review correctly identified that as
\emph{pseudo-replicated} --- generations collapse to 1--2 unique behavioral
fingerprints per (model, arm), so pooling $5\times28$ windows counts the same
underlying defect up to five times and inflates any $\chi^2$. The analysis of
record is therefore at the generation level: we report per-arm uniqueness
counts and an exact two-sided permutation test with the \emph{generation} as
the unit (statistic: difference in mean per-generation pass rate; all
$\binom{10}{5}=252$ relabelings; the smallest attainable $p$ is
$2/252\approx0.0079$, so starred results sit at the test's floor). The
pre-registered checkability policy (conditional vs.\ unconditional scoring)
proved moot: \emph{zero} windows were unscoreable in any constrained or
plain-JS generation.

\subsection{Results}

\begin{table}[ht]
\centering
\caption{Cross-language Phase-3 conformance (mean pass rate over $N{=}5$
generations per cell; 28-window corpus). Conditional and unconditional rates
coincide --- no unscoreable windows.}
\label{tab:crosslang}
\small
\begin{tabular}{lcccc}
\toprule
Model & JS (deployed) & JS (constrained) & plain-JS & \tlaplus{} (constrained) \\
\midrule
Claude Opus 4.8   & 81.4\% & 89.3\%          & \textbf{100\%} & \textbf{100\%} \\
Claude Fable 5    & 57.9\% & 95.7\%          & \textbf{100\%} & \textbf{100\%} \\
Claude Sonnet 4.6 & 50.0\% & \textbf{100\%}  & \textbf{100\%} & \textbf{100\%} \\
Claude Haiku 4.5  & 28.6\% & 40.7\%          & \textbf{100\%} & \textbf{100\%} \\
\bottomrule
\end{tabular}
\end{table}

\begin{table}[ht]
\centering
\caption{Generation-level analysis (replaces a pooled per-window McNemar
table that review identified as pseudo-replicated; those $\chi^2$ values are
retracted). ``Unique'' = distinct behavioral fingerprints (per-window outcome
vectors) among the 5 generations of each arm, ordered deployed-JS /
constrained-JS / plain-JS / \tlaplus{}. $\Delta$ = difference in mean
per-generation pass rate vs.\ \tlaplus{}; $p$ from the exact permutation test
($^{*}$ = $p<0.05$; $0.0079$ is the smallest attainable value). The plain-JS
arm is omitted: it ties \tlaplus{} in every generation ($\Delta{=}0$, $p{=}1$).}
\label{tab:mcnemar}
\small
\begin{tabular}{lccc}
\toprule
Model & Unique (dep/con/plain/\tlaplus{}) & JS(deployed) vs.\ TLA & JS(constrained) vs.\ TLA \\
\midrule
Opus 4.8   & 2/1/1/1 & $\Delta{=}.186$, $p{=}.0079^{*}$ & $\Delta{=}.107$, $p{=}.0079^{*}$ \\
Fable 5    & 2/2/1/1 & $\Delta{=}.421$, $p{=}.0079^{*}$ & $\Delta{=}.043$, $p{=}.44$ \\
Sonnet 4.6 & 1/1/1/1 & $\Delta{=}.500$, $p{=}.0079^{*}$ & $\Delta{=}0$, $p{=}1$ \\
Haiku 4.5  & 1/2/1/1 & $\Delta{=}.714$, $p{=}.0079^{*}$ & $\Delta{=}.593$, $p{=}.0079^{*}$ \\
\bottomrule
\end{tabular}
\end{table}

\subsection{Findings}

\paragraph{Most of the apparent \tlaplus{} advantage was prompt
prescriptiveness, not language.}
Adding the same semantics block to the JavaScript prompt --- changing nothing
else --- moved Fable from 57.9\% to 95.7\% and Sonnet from 50.0\% to a
perfect 100\%. Had we reported the two-arm run, the headline would have been
a decisive language reversal; most of that effect was the prompt. We regard
catching this as a result in its own right: \emph{prompt prescriptiveness is
a first-class experimental variable in specification-generation benchmarks},
large enough to masquerade as a language effect.

\paragraph{A real, one-directional residual survives the prompt control.}
With identical semantics in both prompts, \tlaplus{} reaches 100\%
conformance for \emph{every model and every generation}, while constrained
\jssam{} does so only for Sonnet. At the generation level the residual is
significant for Opus and Haiku (both at the permutation floor,
$p{=}.0079$), vanishes for Sonnet, and for Fable is \emph{not} significant
($\Delta{=}.043$, $p{=}.44$) --- the pooled McNemar had starred Fable's
deficit, and that specific claim was pseudo-replication, now withdrawn. The
residual is largest for the weakest model (Haiku: 40.7\% vs.\ 100\%), and
the direction is uniform: descriptively, no \jssam{} specification ever
passed a window that its \tlaplus{} counterpart failed, in any arm. (When
one arm is at 100\% this is partly structural --- see
\S\ref{sec:limitations} --- but the direction is uniform in the deployed arm
as well, where the gap is significant for all four models.)

\paragraph{The plain-JS arm locates the residual: it is the contract, not
the language.}
The fourth arm ties \tlaplus{} at 100\% for every model and every generation
--- identical per-window outcomes in every paired generation-set
($\Delta{=}0$; no test needed). Haiku is the cleanest single datum:
the weakest model, which reaches only 40.7\% through SAM's machinery,
expresses the identical semantics perfectly as a bare transition function
--- in the same language. Three attributions fall out.
\emph{It is not the language}: same JavaScript, different shape, 100\%.
\emph{It is not formality}: \tlaplus{} and plain-JS share the same shape ---
a minimal declarative single-step transition over exactly the observable
state --- and both sit at 100\%; what the three-arm data suggested was
``formal directness'' is really \emph{minimal declarative shape}.
\emph{It is the SAM contract's defect surface}: proposal/acceptor wiring,
intent firing, and the mandatory auxiliary state
(\texttt{threadStatus}, \texttt{callType}) --- precisely where
Experiment 1's \texttt{try\_lock}-from-free and release bugs lived. Every
moving part the contract obliges is transcription surface, and the tax falls
hardest on the weakest model (Haiku pays 59 points; Sonnet pays nothing).

\paragraph{Verdict against the pre-registered outcomes.}
Not the strong ``\jssam{} wins'' (the direction is reversed), and not a tie
at the backend level --- but at the \emph{language} level, a tie is exactly
what the controlled comparison shows: JavaScript expresses the transition
relation as reliably as \tlaplus{} once the specification shape is matched.
The pre-registered outcome list did not anticipate this resolution; we note
for transparency that the plain-JS arm was designed after the three-arm
results were known (though scored blind, by the same mechanical replay, on
the same corpus).

\paragraph{Design implication for \jssam{}.}
The Phase-3 fidelity gap is the cost of the SAM specification contract, not
a language limit --- and the two can be decoupled. The backend can offer a
lean specification shape (\texttt{next(state, action, data)}, the JavaScript
analogue of the \tlaplus{} relation), or score conformance against a plain
transition core over only the observable variables, dropping the mandatory
auxiliary state from the replay contract. SAM's ceremony is not gratuitous
--- the instance, intents, and value domains are exactly what feed the
Phase-2/4 bounded explorer --- but it should be paid where it buys
something: explorable structure for model checking, not conformance replay.
A prototype confirms the decoupling is practical: a $\sim$40-line generic
breadth-first explorer over \texttt{\{init, next\}} runs the entire pipeline
in the same locked-down sandbox, so the SAM machinery is not \emph{required}
by any phase. Two points keep this precise. First, the SAM checker itself
was never the execution problem: it ran Phases 2 and 4 for every \jssam{} arm
in this paper (the full depth-6 intent exploration each time) without failing
on any specification it was given; the tax is on the \emph{authoring} side ---
the instance, intents, acceptors, and auxiliary state the model must
hand-write to make a specification checker-consumable. (``Ran without
failing'' is itself a bounded claim: the metric audit in \S\ref{sec:models}
shows the checker also ran to completion on a specification whose exploration
never leaves its initial state, so a clean Phase-2 run certifies
checkability, not adequacy.) Second, nothing is lost at this scale by
the replacement: the SAM checker and the generic explorer are both bounded
safety explorers (reachable states, invariants, deadlock), equivalent in
checking power on these state spaces. The hybrid path keeps the library
checker unchanged --- the model writes only the pure \texttt{next}, the
harness generates the scaffolding the checker consumes --- relocating the
checker from the specification contract to the tooling, which is where TLC
has always lived on the \tlaplus{} side: no one asks the model to write
TLC. A follow-up study closes the loop with model-generated
specifications: each of the four models generated five lean specifications
($N{=}5$, 20 in total, each saved and evaluated as a single artifact across
all phases), and \textbf{20 of 20 pass every phase} --- loadable, Phase-2
clean, Phase-3 28/28, all three Phase-4 invariants holding. Two caveats
travel with this result: \texttt{spin}'s observable state is tiny (three
reachable states), so the lean explorer's adequacy must be re-validated on
richer tasks; and the 20 generations collapse to roughly nine unique
solutions --- the models converge on the same correct \texttt{next} ---
which is expected for so small a system but limits the effective sample.
Both caveats are resolved by a second task (\S\ref{sec:locksvc}).

\paragraph{Checkability, revisited.}
Zero unscoreable generations in the constrained \tlaplus{} arm confirms
Experiment 3's attribution: the earlier ``1/4 model-checkable'' was entirely
the unbounded-\texttt{CHOOSE} trap plus the configuration surface, both of
which the pinned contract and direct replay remove. The checkability gap is
real but shallow --- a matter of guardrails, not capability.

\paragraph{A free variance check on Experiment 1.}
The JS(deployed) arm replicates Experiment 1's setting at $N{=}5$. The
single-generation scores (89.3, 50.0, 50.0, 21.4) sit near the five-run means
(81.4, 57.9, 50.0, 28.6), and the model ordering is preserved (with Fable's
mean now slightly above Sonnet's). Experiment 1's qualitative conclusions
survive replication; its exact figures move by up to $\sim$8 points, which is
the error bar readers should apply to any single-generation number in this
paper.

\subsection{Generalizing the lean result: a second task}
\label{sec:locksvc}

The lean-contract demonstration carried two caveats: a three-state
observable space, and 20 generations collapsing to $\sim$9 unique
solutions. To resolve both, we built a second full lean task around
\texttt{locksvc}, SysMoBench's PGo-based distributed lock service ---
chosen over the ring-buffer candidate because it runs as an ordinary Go
test with no kernel build, demonstrating in passing that the trace-capture
methodology of \S\ref{sec:traces} is not kernel-specific. The harness runs
PGo's trace-recorder test five times and folds the raw MPCal events into
observable windows over $\{\mathit{holder}, \mathit{waiters}\}$ --- a lock
holder plus a FIFO wait queue --- yielding a 60-window corpus (15 per
action), self-consistent across runs. The fold also repairs a real trace
artifact: PGo's empty vector clocks allow a \texttt{CriticalSection} event
to be logged before its own \texttt{Grant}; the harness reorders each grant
before its critical section, which is causally sound (the critical section
cannot be entered without the grant). The replay and drivers were
generalized to be task-agnostic: per-task action domains and invariants,
and the projection rule in place of \texttt{spin}-specific keys.

The result replicates in full: \textbf{20 of 20 model-generated lean
specifications pass every phase} (5/5 per model), now over a 31-state
observable space with a genuine queue, and with 11/20 unique solutions
(36--48 lines, clustering by model) --- materially different correct queue
implementations, all conforming. Both caveats are thereby resolved: the
clean sweep is not an artifact of a trivial state space, and the models do
not converge on a single memorized solution. Across the two tasks, 40/40
model-generated lean specifications pass every phase.

\section{Discussion}
\label{sec:discussion}

\paragraph{``Looks right'' is not ``is right,'' in any language.}
The four experiments repeatedly produced specifications that passed every
internal-consistency check while being wrong about the real system --- most
vividly Fable's never-releasing lock, which preserved mutual exclusion
\emph{because} it was broken. Phase 3's value is not incremental rigor; it is
the only phase whose ground truth is external to the model. This replicates
SysMoBench's core claim~\cite{sysmobench,sigops2026} in a new language
substrate, which strengthens it: the phenomenon is about LLM modeling, not
about \tlaplus{}.

\paragraph{The familiar path is modeled; the unfamiliar path is not.}
For the three stronger models, every Experiment-1 failure lands on a
less-common code path: the non-blocking \texttt{try\_lock}-from-free
acquisition (all three) and release bookkeeping (two of the three).
Blocking acquisition --- the version of a spinlock every textbook
presents --- was modeled correctly by all but the smallest model. This is
the ``textbook template'' failure mode of \cite{sigops2026} reproduced at
the granularity of individual actions within a 100-line system.

\paragraph{Repair belongs in the evaluation loop.}
The repair experiment suggests that single-shot leaderboard scores understate
some models (Fable) and overstate the ceiling of others (Haiku). Since
practical formal-modeling deployments will be agentic loops with exactly this
kind of counterexample feedback, benchmarks arguably should report both axes.
The repair driver we contribute makes the second axis cheap to measure for
any \jssam{} task.

\paragraph{The hypothesis, after Experiment 4: a split verdict, sharply
located.}
The practical form of the hypothesis survives: host-language familiarity plus
a self-contained executable artifact gets models to an evaluable
specification far more often (4/4 vs.\ 1/4), and evaluability is the
precondition for everything else. The strong form --- that familiarity yields
more \emph{faithful} models --- resolves to a tie once like is compared with
like: JavaScript in the shape of the \tlaplus{} relation matches \tlaplus{}
at 100\% for every model. What the three-arm data made look like a language
effect was the specification contract: every layer of machinery between the
semantics and the artifact --- proposals, acceptors, auxiliary state --- is
surface on which transcription slips, and the tax falls hardest on the
weakest models. Neither training-data abundance nor formality is the
operative variable on this task; \emph{contract minimality} is. And the
largest single effect in the study is still the prompt
(Sonnet 50\%$\to$100\%): prompt, then contract, then language is the order
of what mattered here.

\section{Limitations and threats to validity}
\label{sec:limitations}

We state these plainly; the study is a case study and should be read as one.

\begin{itemize}
\item \textbf{The controlled comparisons concern one system.} Experiments
  1--4 all concern the spinlock, the \emph{easiest} class of SysMoBench
  task; the lean-contract validation now spans two tasks
  (\S\ref{sec:locksvc}), but the four-arm controlled comparison itself has
  not been repeated on the second. Distributed protocols with richer state
  (Raft variants) may reorder results. The spinlock's simplicity is also
  why Phases 1/2/4 saturate; on harder tasks they discriminate
  too~\cite{sigops2026}.
\item \textbf{Experiments 1--3 are single-generation.} Each cell in
  Tables~\ref{tab:models}--\ref{tab:tladetail} is one sample. Experiment 4's
  $N{=}5$ deployed-JS arm retroactively bounds the damage --- ordering
  preserved, individual scores moving by up to $\sim$8 points --- but the
  repair-loop and checkability experiments have not been replicated and
  should be read with that error bar.
\item \textbf{28 trace windows from two deterministic scenarios.} Contention
  is expressed via \texttt{try\_lock} on a single CPU; true SMP interleavings
  and a blocking waiter's wakeup are not represented. The corpus is faithful
  to real kernel executions but does not exhaust the lock's behavior.
\item \textbf{The \tlaplus{} comparison mixes causes.} The 1/4-vs-4/4 result
  bundles a genuine model-attributable pitfall (unbounded \texttt{CHOOSE},
  2/4) with a harness shortfall (config generation, 1/4). We report the
  attribution per model (Table~\ref{tab:tladetail}); the headline number
  should not be read as purely a language effect. No repair round was run on
  the \tlaplus{} side --- a repaired Opus might fix its \texttt{CHOOSE}
  readily, and symmetry demands that experiment before strong claims.
\item \textbf{Experiment 4 measures transcription, not derivation.} With the
  semantics block present, the constrained and plain arms measure how
  reliably a model \emph{transcribes} near-spec-level semantics into a
  language, not whether it can \emph{derive} the semantics from source code.
  The deployed-JS arm is the closer proxy for derivation --- and there the
  gap is much larger --- but there is no deployed-\tlaplus{} arm to pair it
  with (free-form \tlaplus{} was mostly uncheckable), so the
  derivation-level comparison remains open.
\item \textbf{The plain-JS arm resolves the substrate asymmetry, but was
  added post hoc and answers a narrower question.} The three-arm design left
  ``executable machinery admits bugs a two-variable relation cannot'' as an
  open alternative; the fourth arm was designed to test exactly that, after
  the three-arm results were known (its data were then collected by the same
  blind mechanical replay). Its conclusion has since been extended rather
  than bounded: a generic \texttt{\{init, next\}} explorer replaces the SAM
  scaffolding for Phases 2 and 4 (at equivalent bounded-safety checking
  power on these state spaces), and 20/20 model-generated lean
  specifications pass every phase on \texttt{spin} --- but that
  demonstration inherits \texttt{spin}'s three-state observable space and a
  converged solution set ($\sim$9 unique among 20), so the finding needed a
  second task. The \texttt{locksvc} study (\S\ref{sec:locksvc}) supplies it
  --- a 31-state space, 11/20 unique solutions, 20/20 passing ---
  generalizing the contract-tax removal across tasks of this scale, though
  not yet to large distributed protocols.
\item \textbf{$b{=}0$ is partly structural, and the effective sample is
  small.} When the \tlaplus{} arm passes every window it fails none, so a
  JS-only pass is impossible by construction; any paired comparison reduces
  to whether the \jssam{} failure count is significant. Moreover the
  generations collapse to 1--2 unique behavioral fingerprints per (model,
  arm) --- the pooled window-level McNemar we originally reported counted the
  same defect up to five times, and its $\chi^2$ values are retracted in
  favor of the generation-level permutation test of
  Table~\ref{tab:mcnemar}. At the strictest unit (the unique solution,
  $n{=}1$--$2$ per cell) no significance test is meaningful; the
  per-solution effect sizes stand on their own.
\item \textbf{The direct TLC replay is the controlled condition, not the
  deployed one.} The benchmark ships an agent-mediated \tlaplus{} Phase 3;
  pinning spec variables to the trace schema trades realism for
  comparability. Whether a ``pinned-schema'' Phase-3 mode should exist as a
  first-class controlled condition alongside the agent path is a
  methodological question we are putting to the maintainers.
\item \textbf{Phase 4 differs in template count} (seven \tlaplus{} invariants
  vs.\ three \jssam{}), so it is compared pass/fail only, and the \jssam{}
  invariant set is the weaker filter.
\item \textbf{Model family.} All four models are from one vendor. The
  benchmark's configuration already includes other vendors' models; the
  comparison should be broadened before generalizing.
\item \textbf{Author positionality.} The first author is the author of the
  SAM pattern. The evaluation pipeline is automated and the corpus is
  model-independent, which limits (but does not eliminate) the room for
  favorable design choices; all code, prompts, traces, and drivers are
  public for scrutiny. We note that the study's sharpest negative finding
  --- the contract-tax attribution --- cuts \emph{against} the SAM
  contract as deployed, which we take as evidence the pipeline is not
  tilted toward it.
\end{itemize}

\section{Related work}
\label{sec:related}

\paragraph{Benchmarking LLM formal modeling.}
SysMoBench~\cite{sysmobench,sigops2026} is the direct foundation of this
work: it contributes the four-phase methodology, the eleven-system task
suite, and the finding that conformance --- not syntax --- is where LLM
specifications fail. Our contribution is orthogonal: holding the benchmark
fixed and varying the specification \emph{language}, including the first
non-formal one. The Specula agent~\cite{specula} shows that agentic
workflows can saturate current SysMoBench tasks in \tlaplus{}; our repair
experiment probes the smallest unit of such a workflow (one counterexample
round) across models.

\paragraph{Trace conformance for specifications.}
Validating specifications against implementation traces has industrial
precedent: MongoDB's eXtreme modeling checked \tlaplus{} models against
driver traces~\cite{davis2020extreme}, and trace validation is the standard
answer to ``does the spec match the code''~\cite{sysmobench}. SysMoBench's
transition-window formulation, which we adopt unchanged, makes conformance
per-action and per-window rather than whole-trace.

\paragraph{Formal methods practice.}
\tlaplus{}~\cite{lamport2002specifying} with TLC~\cite{yu1999tlc},
Alloy~\cite{jackson2002alloy}, and CSP-based PAT~\cite{sun2009pat} are the
benchmark's formal backends; industrial deployments
(e.g.,~\cite{newcombe2015amazon}) motivate the whole enterprise. SAM
\cite{sam,dubray2016infoq} occupies an unusual position: an engineering
pattern with \tlaplus{}-derived semantics, here repurposed as a
specification language precisely because it sits at the intersection of
``formally disciplined'' and ``abundant in training data.''

\section{Conclusion and future work}
\label{sec:conclusion}

We extended SysMoBench with \jssam{}, its first non-formal specification
backend, built the kernel trace-capture pipeline needed to give it a real
conformance phase, and ran four experiments on the Asterinas spinlock task,
extending the lean-contract validation to a second task (a distributed
lock service).
The results make four points with unusual precision for a single case study:
transition validation against real traces is the only phase that
discriminated among four frontier models (21.4\%--89.3\%, with all other
phases unanimous); one round of trace-feedback repair reveals a
self-correction axis that one-shot scores conceal --- capable models repair
fully and verifiably (a held-out audit found root-cause fixes, not memorized
windows), the weakest not at all; the same
models that reach an evaluable specification in JavaScript 4/4 times manage
it in \tlaplus{} 1/4 times, for two identifiable structural reasons; and
under a pre-registered, paired, four-arm controlled comparison, the
faithfulness question resolves in an unexpected place --- most of an
apparent \tlaplus{} conformance advantage was prompt prescriptiveness, and
the residual that survives is attributable to the SAM contract's machinery,
not to JavaScript and not to formality: plain JavaScript in the shape of the
\tlaplus{} relation ties \tlaplus{} at 100\% for every model.

The hypothesis that motivated the backend is neither confirmed nor dead; it
is answered more usefully than either. Stated plainly, the study ranks what
governs LLM specification quality: \emph{prompt, then contract shape, then
language}. Familiarity buys \emph{checkability}; contract minimality buys
\emph{transcription fidelity}; and the prompt buys more than either. The
largest effect anywhere in the study came from four lines of prompt
(Sonnet, 50\%$\to$100\%). The next came from the shape of the artifact: a
minimal declarative transition function over exactly the observable state
was written flawlessly by every model, even the weakest, in either language
--- while every layer of ceremony between the semantics and the artifact
cost conformance. The language itself, once shape was matched, contributed
nothing measurable. On the practical question --- can JavaScript replace
\tlaplus{}? --- the evidence splits cleanly. For LLM-generated conformance
modeling at this scale: yes, interchangeably on fidelity, and with fewer
operational failure modes (no \texttt{CHOOSE} trap, no configuration
artifact). For deep verification: no --- TLC's temporal properties,
liveness checking, and large-state model checking have no equivalent in
either bounded safety explorer used here (the SAM library's checker or its
40-line generic replacement), and remain \tlaplus{}'s distinctive value. The search
this motivates is therefore not for a better specification \emph{language}
but for a better specification \emph{contract}:
\texttt{\{init, next\}} over observable keys, in whatever language the
model writes well, with richer machinery generated by the harness rather
than transcribed by the model. Two questions remain genuinely open:
whether these conclusions survive \emph{derivation} --- models extracting
the semantics from source code too complex to guess --- and whether they
survive scale.

The experiments that follow are concrete: (1)~promoting the lean
specification shape from demonstration to backend: with model-generated
lean specifications passing all phases 40/40 across both tasks, the
remaining question is purely integration --- a first-class lean backend, a
hybrid where the model writes only the pure core and the harness generates
the SAM scaffolding, or the documented status quo; (2)~a derivation-level
comparison --- both languages, \emph{no} semantics block, on a task whose
semantics cannot be guessed from the action names --- now unblocked by the
\texttt{locksvc} harness and still the decisive open test of the
familiarity claim; (3)~a prompt-prescriptiveness
dose--response study, treating the semantics block as the manipulated
variable rather than a nuisance; (4)~an as-deployed agent-path arm beside
the controlled replay, for ecological validity; (5)~a \tlaplus{} repair
round for symmetry with Experiment 2; and (6)~replication beyond one model
family. We are contributing the backend, harnesses, drivers, corpora, and
the controlled-comparison methodology upstream so these can be run --- and
challenged --- by the community.

\section*{Acknowledgments}
We thank the SysMoBench / Specula maintainers for the benchmark, the
integration guidance that shaped the \jssam{} backend's scope, and the
reference instrumentation for the Asterinas spinlock.

\section*{Reproducibility}
All artifacts are public: the \jssam{} backend
(\texttt{tla\_eval/languages/js\_sam.py}, \texttt{tools/js-sam/}), the trace
harness (\texttt{scripts/harness/spin/}), the captured corpus
(\texttt{data/sys\_traces/spin/}), the repair driver
(\texttt{scripts/repair\_phase3.py}), the Experiment-4 apparatus --- the
constrained and plain-JS prompts, the direct TLC replay and its
functional-check audit (\texttt{scripts/tla\_direct\_tv.py}, the 20 audited
specifications in \texttt{output/tla\_specs/},
\texttt{output/tla\_functional\_audit.json}, the permissive control in
\texttt{tools/tla/}), the paired four-arm driver with raw
per-window data (\texttt{scripts/tla\_phase3\_study.py},
\texttt{output/tla\_phase3\_study.json}), the pre-registered design and
results (\texttt{docs/js\_sam\_tla\_phase3\_*.md}), the lean-contract
prototype and full-pipeline studies (\texttt{tools/plain-js/},
\texttt{scripts/plain\_js\_tv.py},
\texttt{scripts/lean\_full\_pipeline\_study.py}, the 40 saved generations in
\texttt{output/lean\_specs*/}, \texttt{docs/js\_sam\_lean\_contract.md}),
the \texttt{locksvc} harness and corpus
(\texttt{scripts/harness/locksvc/}, \texttt{data/sys\_traces/locksvc/}), and
the methodology note
(\texttt{paper/methodology.md}) --- and the experiment write-ups
(\texttt{docs/js\_sam\_*.md}) in the SysMoBench fork at
\url{https://github.com/jdubray/SysMoBench-1}, staged as stacked pull
requests against \url{https://github.com/specula-org/SysMoBench}.
Requirements: Docker, Node $\geq$ 20, Python 3, and (for the \tlaplus{} arm)
Java with \texttt{tla2tools}. A single script reproduces the kernel capture:
\texttt{scripts/harness/spin/run.sh}.

\begin{thebibliography}{10}
\small

\bibitem{sysmobench}
Q.~Cheng, R.~Tang, E.~Ma, F.~Hackett, P.~He, Y.~Su, I.~Beschastnikh,
Y.~Huang, X.~Ma, and T.~Xu.
\newblock SysMoBench: Evaluating AI on formally modeling real-world systems.
\newblock arXiv:2509.23130, 2025. Leaderboard: \url{https://sysmobench.com}.

\bibitem{sigops2026}
Q.~Cheng, R.~Tang, E.~Ma, F.~Hackett, P.~He, Y.~Su, I.~Beschastnikh,
Y.~Huang, X.~Ma, and T.~Xu.
\newblock Can LLMs model real-world systems in TLA+?
\newblock ACM SIGOPS Blog, May 2026.
\newblock \url{https://www.sigops.org/2026/can-llms-model-real-world-systems-in-tla/}.

\bibitem{lamport2002specifying}
L.~Lamport.
\newblock \emph{Specifying Systems: The TLA+ Language and Tools for Hardware
  and Software Engineers}.
\newblock Addison-Wesley, 2002.

\bibitem{yu1999tlc}
Y.~Yu, P.~Manolios, and L.~Lamport.
\newblock Model checking TLA+ specifications.
\newblock In \emph{CHARME}, 1999.

\bibitem{jackson2002alloy}
D.~Jackson.
\newblock Alloy: A lightweight object modelling notation.
\newblock \emph{ACM TOSEM}, 11(2), 2002.

\bibitem{sun2009pat}
J.~Sun, Y.~Liu, J.~S. Dong, and J.~Pang.
\newblock PAT: Towards flexible verification under fairness.
\newblock In \emph{CAV}, 2009.

\bibitem{newcombe2015amazon}
C.~Newcombe, T.~Rath, F.~Zhang, B.~Munteanu, M.~Brooker, and M.~Deardeuff.
\newblock How Amazon Web Services uses formal methods.
\newblock \emph{Communications of the ACM}, 58(4), 2015.

\bibitem{davis2020extreme}
A.~J. Davis, M.~Hirschhorn, and J.~Schvimer.
\newblock eXtreme modelling in practice.
\newblock \emph{PVLDB}, 13(9), 2020.

\bibitem{sam}
J.-J. Dubray.
\newblock The SAM pattern.
\newblock \url{https://sam.js.org}.

\bibitem{dubray2016infoq}
J.-J. Dubray.
\newblock Why I no longer use MVC frameworks.
\newblock InfoQ, February 2016.
\newblock \url{https://www.infoq.com/articles/no-more-mvc-frameworks/}.

\bibitem{sampattern}
J.-J. Dubray.
\newblock \texttt{sam-pattern}: a library implementing the SAM pattern.
\newblock \url{https://www.npmjs.com/package/@cognitive-fab/sam-pattern}.

\bibitem{asterinas}
Asterinas project.
\newblock Asterinas: a Linux ABI-compatible, Rust-based framekernel OS.
\newblock \url{https://github.com/asterinas/asterinas}.

\bibitem{specula}
Specula project.
\newblock Specula: an agent specialized in TLA+ formal modeling.
\newblock \url{https://github.com/specula-org/Specula}.

\end{thebibliography}

\end{document}
